% =====================================================================
%  Capítulo 3 — Trabalhos relacionados
% =====================================================================
\chapter{Trabalhos Relacionados}
\label{cap:trabalhos}

Este capítulo situa o trabalho em relação à literatura e a sistemas correlatos.
A revisão aqui apresentada é delimitada (não sistemática) e foca nas três
dimensões do problema: monitoramento de condição por vibração, geração de
sinais sintéticos e conjuntos de dados para treinamento, e estratégias de
armazenamento seletivo de dados de vibração.

\section{Monitoramento de condição e análise de vibração}

O monitoramento de condição por vibração é um campo maduro, consolidado em
normas e em sistemas comerciais. Randall \cite{randall2011} apresenta o estado
da arte das técnicas de monitoramento baseado em vibração para aplicações
industriais, aeroespaciais e automotivas, incluindo a análise de envelope para
falhas incipientes de rolamento e a interpretação de bandas laterais. Scheffer
e Girdhar \cite{scheffer2004} oferecem uma abordagem prática da análise de
vibração e da manutenção preditiva, cobrindo o fluxo completo desde a
aquisição até o diagnóstico e a recomendação de ação. O manual de Harris e
Piersol \cite{harris2002} é referência clássica para as bases teóricas da
dinâmica e da instrumentação de vibração.

Sistemas comerciais de monitoramento (provenientes de fornecedores de
instrumentação e de plataformas de manutenção preditiva) implementam
funcionalidades semelhantes às deste trabalho --- aquisição, FFT, alarmes e
tendências ---, porém tipicamente como produtos fechados, acoplados a hardware
específico de aquisição. O diferencial do ARGUS não está na análise espectral em
si, mas na combinação de \emph{geração sintética de defeitos} e
\emph{persistência seletiva} em um sistema aberto, extensível e validado por
testes, orientado ao treinamento e à pesquisa.

\section{Bases de dados e sinais sintéticos para treinamento}

Para treinamento de analistas e de algoritmos de aprendizagem de máquina, a
comunidade utiliza tanto bases de dados medidas em bancadas de ensaio quanto
sinais sintéticos. As bases medidas, ainda que reais, apresentam limitações
importantes: cobrem um conjunto restrito de máquinas, condições e defeitos;
exigem a reprodução física da falha (com custo e risco); e a documentação dos
estágios de severidade nem sempre é completa. Sinais sintéticos, por sua vez,
permitem gerar grandes volumes rotulados com controle de severidade, de ruído e
de condição operacional, a custo desprezível e sem risco --- exatamente o
propósito do gerador desenvolvido neste trabalho. A limitação inerente é que as
amplitudes sintéticas são parâmetros de síntese e não refletem a dinâmica real
de uma máquina específica (massas, rigidez, amortecimento, ressonâncias), como
discutido no Capítulo~\ref{cap:fundamentacao} e retomado no
Capítulo~\ref{cap:discussao}.

A especificação de assinaturas adotada no projeto segue a literatura de
diagnóstico para definir quais componentes espectrais (ordens e frequências de
falha) devem estar presentes em cada defeito, o que confere coerência física aos
sinais gerados, ainda que as amplitudes sejam relativas.

\section{Persistência seletiva de dados de vibração}

O problema da ``massa de dados'' no monitoramento de condição é reconhecido na
indústria, e a literatura de manutenção preditiva aponta a necessidade de
transformar o dado bruto em informação antes do armazenamento
\cite{scheffer2004}. A abordagem do ARGUS --- representar a leitura por vetores
de picos \(R^3\) e métricas de energia, e persistir apenas transições de estado
segundo um Paralelepípedo de Descarte --- segue essa diretriz e a materializa
em um motor de descarte testável. A comparação contra a última leitura
\emph{efetivamente persistida} (e não contra a última leitura avaliada) é um
detalhe de projeto que evita o acúmulo de \emph{drift}, aspecto que distingue a
implementação de um descarte ingênuo por diferença da leitura anterior.

\section{Lacuna e contribuição deste trabalho}

A partir da revisão delimitada, identifica-se a seguinte lacuna: ferramentas
abertas que integrem, de forma validada e testável, (i) a geração sintética de
sinais de vibração com assinaturas espectrais de um catálogo amplo de defeitos,
(ii) a persistência seletiva segundo o critério \(R^3\) e (iii) uma interface
web para análise e registro de \emph{snapshots} de defeito. O ARGUS propõe-se a
preencher essa lacuna como uma plataforma de treinamento e pesquisa, apoiada em
documentação de requisitos, design e testes que garantem reprodutibilidade.

Uma revisão bibliográfica sistemática e o aprofundamento quantitativo das
comparações com bases de dados públicas são registrados como pendências do
trabalho (ver `MONOGRAFIA_PENDENCIAS.md`), devendo ser conduzidos com o
orientador antes da versão final.
